%قبل از شروع مسئله ترجیحاً چیزی ننویسید. اما اگر خواستید بنویسید، در قالب لیستی (فقرات و شمارش) بنویسید.
%برای اضافه کردن نمره، بعدی از شروع‌مسئله، [x] را اضافه کنید و به جای x  نمره مورد نظر را قرار دهید. برای زیرمسئله‌ها هم به همین صورت است.
%\شروع‌مسئله[10]

\شروع‌مسئله[]
تعریف فرضی: چپ‌ترین نوه را برای یک راس مانند $v$ راسی میگوییم که از به بچه‌ی چپ رفتن متوالی با شروع از راس $v$ به دست می‌آید.

به راحتی می‌توان مشاهده کرد که چپ‌ترین نوه‌ی $v$، کوچکترین راس در زیردرخت به ریشه‌‌ی $v$ می‌باشد.

توضیح الگوریتم روی رأس $v$:
اگر $v$ فرزند راست داشت جواب چپ‌ترین نوه‌ی فرزند راست $v$می‌باشد. چون که فرزندان چپ $v$ از آن کوچکترند و سایر اجزای درخت نیز یا از کل رئوس زیردرخت $v$ کوچکترند یا بزرگتر. 

در صورت برقرار نبودن شرط بالا، با شروع از راس $v$ برای یک رأس در صورتی که فرزند راست پدرش باشد به پدرش می‌رویم. هنگامی که یک رأس فرزند چپ پدرش بود، پدر آن رأس اولین عضو بزرگ بعد $v$ خواهد بود.
اگر روش بالا جواب نداد حتما در بزرگترین رأس درخت بوده‌ایم.

تحلیل زمانی:

فرض کنید که $k$ از اندازه‌ی زیردرخت فرزند راست رأس $v$ نابیشتر باشد. در این حالت الگوریتم مانند یک پیمایش inorder روی فرزند راست رأس $v$ عمل می کند و می‌دانیم که به دست آوردن $k$  عنصر اول inorder از مرتبه‌‌ی
$O(h+k)$
می‌باشد. (به راحتی میتوانید به صورت تفکر بازگشتی و به کم استقرا این موضوع را بررسی نمایید)

حال فرض کنید که $k$ از اندازه‌ی زیردرخت فرزند راست رأس $v$ بیشتر باشد. در این صورت پس از پیمایش inorder کل زیردرخت راست $v$ در مرتبه‌ی زمانی خطی نسبت به اندازه‌ی آن، حالت دوم الگوریتم رخ داده و شروع به رفتن به پدر می‌کنیم تا به پدری برسیم که از فرزند چپ آن وارد شده باشیم. سپس برای این راس ابتدا زیردرخت راست آن‌را (در حالت تهی بودن یک زیردرخت تهی) را می‌پیماییم و دوباره به پدران رئوس می‌رویم، مشاهده می‌شود که رئوسی که فرزند راست آن‌ها باز می‌شود تشکیل یک مسیر از پایین درخت به بالای آن را می‌دهند. و هر راسی که سمت راست آن باز میشود پدر راس‌های قبلی با این ویژگی بود است.

طول این مسیر از $O(h)$ است و تنها یک بار از پایین به بالا پیموده می‌شود، بقیه‌ی هزینه مربوط به پیمایش inorder فرزندان راست رئوس این مسیر می‌باشد که هزینه‌ی آن برابر اندازه‌ی زیردرخت می‌باشد و به اندازه‌ی زیردرخت از $k$ تا مصرف می‌کند.

پس زمان الگوریتم از مرتبه‌ی طول مسیر که حداکثر $O(h)$ و مرتبه‌ی مجموع اندازه‌ی زیر درخت‌ها برای پیمایش inorder آن‌ها به علاوه‌ی یک h برای پیمایش ناقص زیردرخت آخر یعنی از مرتبه‌ی $O(h+k)$ می‌باشد و مجموعا $O(h+k($ می‌باشد.

شکل زیر شهود بیشتری درباره‌ی چگونگی پیمایش مسیر توسط الگوریتم و تحلیل زمانی به شما می‌دهد.

\begin{center}
\درج‌تصویر[پهنا=1.0\پهنای‌متن]{src/Theorical/traverse.png}
\end{center}


 

%اگر مسئله چند قسمت داشت به این صورت عمل کنید:
%\شروع{زیرمسائل}
%\زیرمسئله[0] % مسئلۀ ۱آ
%\زیرمسئله % مسئلۀ ۱ب
%\زیرمسئله % مسئلۀ ۱پ
%\زیرمسئله % مسئلۀ ۱ت
%\پایان{زیرمسائل}

\solution{
	پاسخ نمونه.
}




%%%%%%%%%%%%%%%%%%%%%%%%%%%%%%%%%%%%%%%%%%%%%%%%%%%%%
% برخی دستورات مفید برای تمرین‌ها. %
%%%%%%%%%%%%%%%%%%%%%%%%%%%%%%%%%%%%%%%%%%%%%%%%%%%%%

% فرمول‌های نمونه:
%$f(x) = \Theta(x)$
%$T(x, y) \leq \log(x) + 2^y + \binom{2n}{n}$
% {\tt code\_function}


% لیست موارد بدون شماره‌گذاری:
%\شروع{فقرات}
%    \فقره مورد اول 
%        \شروع{فقرات}
%            \فقره مورد اول داخلی 
%                \شروع{فقرات}
%                    \فقره اولین مورد خیلی داخلی 
%                    \فقره دومین مورد خیلی داخلی 
%                \پایان{فقرات}
%            \فقره مورد دوم داخلی 
%        \پایان{فقرات}
%    \فقره مورد دوم 
%\پایان{فقرات}

% لیست موارد شماره‌گذاری شده:
%\شروع{شمارش}
%    \فقره یک مورد 
%    \فقره یک مورد دیگر 
%    \فقره باز هم یک مورد دیگر
%\پایان{شمارش}

% برای نوشتن یک فرمول چند خطی که یک شماره بخوره. برای حالت‌های دیگر به سایت زیر مراجعه کنید:
% https://www.overleaf.com/learn/latex/Aligning_equations_with_amsmath
%\begin{align} 
%    \begin{split}
%        (x+y)^3 &= (x+y)^2(x+y) \\
%                &= (x^2+2xy+y^2)(x+y) \\
%                &= (x^3+2x^2y+xy^2) + (x^2y+2xy^2+y^3) \\
%                &= x^3+3x^2y+3xy^2+y^3
%    \end{split}                                 
%\end{align}

% نوشتن ماتریس و آرایه دوبعدی
%\begin{align}
%    A = 
%    \begin{bmatrix}
%        A_{11} & A_{21} \\
%        A_{21} & A_{22}
%    \end{bmatrix}
%\end{align}

% اضافه کردن تصویر با شرح. برای اضافه کردن تصویر بدون شرح، فقط دستور درج‌تصویر را استفاده کنید.
%\شروع{شکل}[h!]
%\تنظیم‌ازوسط
%  \درج‌تصویر[پهنا=0.5\پهنای‌متن]{img.jpg}
%\شرح{شرح تصویر}
%\برچسب{تصویر:۱}
%\پایان{شکل}
%\پایان{زیرمسائل}

%کد زیر را ببینید.
%%\lstinputlisting{source_filename.py}
%\begin{latin}
%\begin{lstlisting}
%# Import the modules
%import sys
%import random
%
%ans = True
%
%while ans:
%question = raw_input("Ask the magic 8 ball a question: (press enter to quit) ")
%
%answers = random.randint(1,8)
%
%if question == "":
%sys.exit()
%
%elif answers == 1:
%print "It is certain"
%
%elif answers == 2:
%print "Outlook good"
%
%elif answers == 3:
%print "You may rely on it"
%
%elif answers == 4:
%print "Ask again later"
%
%elif answers == 5:
%print "Concentrate and ask again"
%
%elif answers == 6:
%print "Reply hazy, try again"
%
%elif answers == 7:
%print "My reply is no"
%
%elif answers == 8:
%print "My sources say no"
%\end{lstlisting}
%\end{latin}